% this replaces the keystroke symbols from the keystroke package
\usepackage[os=win]{menukeys}

% create a macro for keyboard keys
\newmenumacro\key{roundedkeys}

% map it to \keystroke commands so we don't need to change the sources after migrating from
% the keystroke package
\newcommand{\keystroke}[1]{\raisebox{-1pt}{\key{#1}}}

% redefine the commands for common keys
\newcommand{\Ctrl}{\keystroke{\ctrl}}
\let\alt\Alt
\renewcommand{\Alt}{\keystroke{\alt}}
\let\oldshift\shift
\renewcommand{\shift}{Shift\,\oldshift}
\newcommand{\Shift}{\keystroke{\shift}}
\newcommand{\Home}{\keystroke{Home}}
\newcommand{\UArrow}{\keystroke{\arrowkeyup}}
\newcommand{\DArrow}{\keystroke{\arrowkeydown}}
\newcommand{\RArrow}{\keystroke{\arrowkeyright}}
\newcommand{\LArrow}{\keystroke{\arrowkeyleft}}
\newcommand{\Enter}{\keystroke{Enter}}
\newcommand{\Spacebar}{\keystroke{Space}}

% prepend a key with Shift (encoding for MFD button shortcuts)
\newcommand{\mfdkey}[1]{\keystroke{\shift,{#1}}}

% patch the roundedkeys style so that keystroke symbols appearing in
% two consecutive lines of running text don't stretch the line spacing
\makeatletter
\tikzset{
  tw@roundedkeys@base/.append style={
    text height=1.74ex,     % Modify the inner height spacing
    text depth=0.67ex       % Modify the inner depth spacing
  }
}
\makeatother

% patch to replace the '+' key separator used for comma-separated
% arguments with a thin space
\makeatletter
\tw@declare@style@simple*{roundedkeys}{%
  \tikz[baseline={($(tw@node.base)+(0,-0.2ex)$)}]{%
    \node(tw@node)[tw@roundedkeys@base]%
    {\strut\color{\usemenucolor{txt}}\CurrentMenuElement};}%
}[%
  \kern 0.08em
]{gray}
\makeatother

% Now for MFD buttons. First, define a colour theme to distinguish
% from keyboard key symbols
\newmenucolortheme{mfdcol}{rgb}{0.97,1,0.97}{0.3,0.3,0.3}{0,0,0}

% Define a new menukeys style to represent MFD button labels
% Arguments for \newmenustyle are:
% {name}[pre]{first}[sep]{mid}{last}{single}[post]{theme}
\newmenustyle{monofixedkeys}{%
  % TikZ styling applied to the key node
  draw=\usemenucolor{br},
  fill=\usemenucolor{bg},
  rounded corners=1pt,
  font=\relsize{-0.8}\ttfamily,
  % Force a minimum width to cleanly fit 3 monospaced characters
  minimum width=2.2em, 
  align=center,
  % Maintain a uniform box height and depth regardless of descenders
  text height=1.74ex,
  text depth=0.67ex,
  inner sep=0pt,
  inner xsep=2pt
}{%
  % Style for sequences: reuse the single style for mid, last, single elements
  draw=\usemenucolor{br},
  fill=\usemenucolor{bg},
  rounded corners=1pt,
  font=\relsize{-0.8}\ttfamily,
  minimum width=2.2em,
  align=center,
  text height=1.74ex,
  text depth=0.67ex,
  inner sep=0pt,
  inner xsep=2pt
}{%
  draw=\usemenucolor{br},
  fill=\usemenucolor{bg},
  rounded corners=1pt,
  font=\relsize{-0.8}\ttfamily,
  minimum width=2.2em,
  align=center,
  text height=1.74ex,
  text depth=0.67ex,
  inner sep=0pt,
  inner xsep=2pt
}{%
  draw=\usemenucolor{br},
  fill=\usemenucolor{bg},
  rounded corners=1pt,
  font=\relsize{-0.8}\ttfamily,
  minimum width=2.2em,
  align=center,
  text height=1.74ex,
  text depth=0.67ex,
  inner sep=0pt,
  inner xsep=2pt
}{gray}
% Bind the style to the mfdbtn macro

% create the new macro for the MFD buttons from the style defined above
\newmenumacro{\mfdbtn}{monofixedkeys}

% and apply the colour theme
\changemenucolortheme{monofixedkeys}{mfdcol}
